\documentclass[11pt,a4paper]{article}
\usepackage[T1]{fontenc}
\usepackage{isabelle,isabellesym}
\usepackage{amsmath,amssymb}
\usepackage{pdfsetup}

\urlstyle{rm}
\isabellestyle{it}

\begin{document}

\title{K\H{o}nig's Edge Colouring Theorem}
\author{Ari Gilchrist}
\maketitle

\begin{abstract}
  We formalise K\H{o}nig's edge colouring theorem (1916): for any finite bipartite graph, the
  chromatic index is the maximum degree. The proof follows the classical
  Kempe chain argument, with the chain defined in terms of alternating
  walks rather than maximal simple paths. As a corollary, we show that the edge set of every finite $k$-regular bipartite graph is the disjoint
  union of $k$ perfect matchings. Our development builds on the
  \emph{Undirected Graph Theory} library.
\end{abstract}

\subsection*{Acknowledgements}
This formalisation was generated almost entirely by Anthropic's Claude
Opus 5.5, which designed the proof strategy, wrote and debugged the Isabelle code, and
drafted the accompanying document. The author only reviewed the development,
revised the text, and verified that the definitions and theorem statements capture the intended
mathematics.

\tableofcontents

\parindent 0pt\parskip 0.5ex

\section{Introduction}

A \emph{proper edge colouring} of a graph assigns a colour to every edge such
that any two edges sharing a vertex receive different colours. The
\emph{chromatic index} $\chi'(G)$ of a graph $G$ is the minimum number of
colours for which such a colouring exists. Since all edges connected to a vertex of
maximum degree $\Delta(G)$ must receive pairwise distinct colours, every graph trivially
satisfies $\chi'(G) \geq \Delta(G)$. In the special case of bipartite graphs,
K\H{o}nig~\cite{koenig1916} proved that the lower bound is attained:

\begin{quote}
  \textbf{Theorem (K\H{o}nig, 1916).} For every finite bipartite graph $G$,
  $\chi'(G) = \Delta(G)$.
\end{quote}

The substantial direction is therefore $\chi'(G) \leq \Delta(G)$, i.e.\ the existence
of a proper edge colouring with $\Delta(G)$ colours. To find one, we follow the standard
proof~\cite[Proposition~5.3.1]{diestel}, which considers colouring each edge at a time.
At the moment that a new edge $\{x,y\}$ is considered for colouring, the vertices $x$
and $y$ have at most
$\Delta - 1$ coloured edges, hence some colour $\alpha$ is missing at $x$ and
some colour $\beta$ is missing at $y$. If $\alpha$ is also missing at $y$, the
new edge is coloured $\alpha$, otherwise the colours $\alpha$ and $\beta$ are swapped
along the $\alpha/\beta$-alternating chain (or \emph{Kempe chain}) starting
at $y$. Bipartiteness
ensures that this chain cannot reach $x$. After the swap, $\alpha$ is
therefore missing at both endpoints, and thus $\{x,y\}$ may be assigned the colour $\alpha$
while preserving the colouring's properness.

Rather than reasoning about maximal simple alternating paths, we define the
chain as the set of all edges traversed by alternating walks from $y$.
The key observation is a closure property: at every vertex, either all or
none of its $\alpha$- and $\beta$-coloured edges belong to the chain.
Consequently, swapping the two colours on the chain preserves properness.

Section~2 defines proper edge colourings, develops the
alternating chain argument, and proves K\H{o}nig's theorem.
Section~3 derives K\H{o}nig's original application:
the edge set of every finite $k$-regular bipartite graph decomposes into $k$
perfect matchings. The development builds on the \emph{Undirected Graph
Theory} entry~\cite{Undirected_Graph_Theory-AFP} of the Archive of Formal Proofs.
Following that library, a bipartite graph comes with a partition of its
vertices into two non-empty parts. This excludes only graphs with fewer than
two vertices, which have no edges.

\input{session}

\bibliographystyle{abbrv}
\bibliography{root}

\end{document}
